\begin{proposition}[3.4.6]
\label{IV.2.3.4.6-fr}
Siendo las hipótesis las de \textup{(3.4.1)}, sea $x$ un punto de $Y$. Supongamos
que $Y$ no contiene ninguno de los $S_i$ que contienen a $x$, y que
$\mathcal F/t\mathcal F$ es reducido en el punto $x$ \textup{(3.2.2)}. Entonces $t_x$ es
$\mathcal F_x$-regular y $\mathcal F$ es reducido en el punto $x$. Además, si
$z_j$ es el punto genérico de una componente irreducible de
$\operatorname{Supp}(\mathcal F/t\mathcal F)$ que contiene a $x$, $z_j$ está
contenido en uno solo de los $S_i$, y $\mathcal O_{S,z_j}$ es un anillo de
valoración discreta cuya uniformizante es $t_{z_j}$.
\end{proposition}

\begin{proof}
El hecho de que $t_x$ sea $\mathcal F_x$-regular resulta del lema siguiente,
aplicado al anillo $\mathcal O_x$:

\begin{lemma}[3.4.6.1]
\label{IV.2.3.4.6.1-fr}
Sean $A$ un anillo noetheriano, $M$ un $A$-módulo de tipo finito,
$\mathfrak p_i$ los elementos minimales de $\operatorname{Supp}(M)$, $t$ un
elemento de $A$. Se supone que $t$ no pertenece a ninguno de los
$\mathfrak p_i$ y que $M/tM$ es un $A$-módulo reducido \textup{(3.2.2)}. Entonces $t$
es $M$-regular.
\end{lemma}

\begin{proof}
Todo ideal primo $\mathfrak p\in\operatorname{Supp}(M)$ contiene uno de los
$\mathfrak p_i$; como $t$ no pertenece a ninguno de los $\mathfrak p_i$, la
homotecia de razón $t$ en $M_{\mathfrak p}$ no es nilpotente (Bourbaki,
\emph{Alg. comm.}, cap.~IV, \S~1, núm.~4, cor. de la prop.~9). Denotemos por
$N$ el submódulo de $M$ formado por los elementos anulados por una potencia de
$t$, y pongamos $P=M/N$; mostraremos que $N=0$. Como $t$ es $P$-regular, se
tiene la sucesión exacta (3.4.1.4)
\[
0\longrightarrow N/tN\longrightarrow M/tM\longrightarrow P/tP
\longrightarrow0.
\]
Como $N$ es de tipo finito, está anulado por una potencia de $t$, y basta
demostrar que $N/tN=0$. Como $N/tN$ es un submódulo de $M/tM$, basta probar
que $(N/tN)_{\mathfrak p}=0$ para todo
$\mathfrak p\in\operatorname{Ass}(M/tM)$, o, equivalentemente, que el
homomorfismo
$u_{\mathfrak p}:(M/tM)_{\mathfrak p}\longrightarrow(P/tP)_{\mathfrak p}$
es biyectivo para todo $\mathfrak p\in\operatorname{Ass}(M/tM)$. Ahora bien,
$(P/tP)_{\mathfrak p}\neq0$; en efecto, como
$\mathfrak p\in\operatorname{Supp}(M/tM)
=\operatorname{Supp}(M)\cap V(t)$ $(\mathbf{0}_{\mathrm I},1.7.5)$, la imagen de $t$
en $A_{\mathfrak p}$ está contenida en el ideal maximal
$\mathfrak pA_{\mathfrak p}$, por lo que la hipótesis
$P_{\mathfrak p}=tP_{\mathfrak p}$ implicaría $P_{\mathfrak p}=0$ por el lema
de Nakayama; se tendría entonces $M_{\mathfrak p}=N_{\mathfrak p}$ y la
homotecia de razón $t$ en $M_{\mathfrak p}$ sería nilpotente; pero esto
contradice la observación hecha al principio, puesto que
$\mathfrak p\in\operatorname{Supp}(M)$. Dicho esto, la hipótesis de que $M/tM$
es reducido implica $\operatorname{long}((M/tM)_{\mathfrak p})=1$, y como
$(P/tP)_{\mathfrak p}\neq0$, $u_{\mathfrak p}$ es necesariamente biyectivo,
lo que demuestra el lema.
\end{proof}

Por hipótesis, ninguno de los ciclos primos asociados inmersos de
$\mathcal F/t\mathcal F$ contiene a $x$, por lo que ninguno de los ciclos
primos asociados inmersos de $\mathcal F$ contiene a $x$, en virtud de
(3.4.4). Por otra parte, aplicando (3.4.2) a una componente irreducible de
$\operatorname{Supp}(\mathcal F/t\mathcal F)$ que contiene a $x$, se ve que
$\operatorname{long}(\mathcal F_{s_i})=1$ para todo $S_i$ que contiene a $x$,
lo que termina de demostrar que $\mathcal F_x$ es reducido; finalmente, las
últimas afirmaciones son también consecuencias de (3.4.2).
\end{proof}

\begin{corollary}[3.4.7]
\label{IV.2.3.4.7-fr}
Sean $A$ un anillo local noetheriano, $\mathfrak m$ su ideal maximal, $M$ un
$A$-módulo de tipo finito, $(x_i)_{1\leqslant i\leqslant k}$ una familia de
elementos de $\mathfrak m$ que forma parte de un sistema de

\oldpage[IV-2]{51}

parámetros para $M$ $(\mathbf{0},16.3.6)$. Si el $A$-módulo
$N=M/(\sum_{i=1}^k x_iM)$
es íntegro \textup{(3.2.4)}, entonces $M$ es íntegro y la sucesión
$(x_i)_{1\leqslant i\leqslant k}$ es $M$-regular.
\end{corollary}

\begin{proof}
Por inducción sobre $k$, la cuestión se reduce inmediatamente al caso $k=1$; escribiremos
$x$ en lugar de $x_1$; la hipótesis de que $x$ forma parte de un sistema de
parámetros para $M$ implica que
$\dim(N)=\dim(M)\setmin 1$ $(\mathbf{0},16.3.7)$. Pongamos $n=\dim(N)$; existe, por tanto,
un ideal primo minimal $\mathfrak p$ de $\operatorname{Supp}(M)$ tal que
$\dim(M/\mathfrak pM)=n+1$ $(\mathbf{0},16.3.4)$, y para todo entero $j>0$ también se
tiene $\dim(M/\mathfrak p^jM)=n+1$ $(\mathbf{0},16.3.5)$; además, $x$ forma parte de
un sistema de parámetros para $M/\mathfrak p^jM$ $(\mathbf{0},16.3.5)$, por lo que,
si se pone $M'=M/\mathfrak p^jM$ y $N'=M'/xM'$, se tiene $\dim(N')=n$. Es
claro que existe un homomorfismo sobreyectivo $v:N\to N'$; demostremos que
$v$ es \emph{biyectivo}. En efecto, si $P=\operatorname{Ker}(v)$, se tiene
$\operatorname{Ass}(P)\subset\operatorname{Ass}(N)$ y, puesto que $N$ es
íntegro, la hipótesis $P\neq0$ implicaría que tanto $\operatorname{Ass}(P)$
como $\operatorname{Ass}(N)$ se reducirían al único punto
$\mathfrak q$ de $\operatorname{Ass}(N)=\operatorname{Supp}(N)$; pero como
$\dim(N')=\dim(N)$, se tiene $N'_{\mathfrak q}\neq0$, y la hipótesis
$\operatorname{long}(N_{\mathfrak q})=1$ implica
$\operatorname{long}(N'_{\mathfrak q})=1$, ya que
$N_{\mathfrak q}\to N'_{\mathfrak q}$ es sobreyectivo. Se tendría entonces
$P_{\mathfrak q}=0$, en contradicción con la hipótesis, de donde se sigue
nuestra afirmación. Pero entonces $N'$, al ser isomorfo a $N$, es íntegro;
además, el soporte de $N'$ (igual a la intersección de
$\operatorname{Supp}(M)$ y $V(x)$) no puede contener
$\operatorname{Supp}(M')$, y este último conjunto es irreducible por
construcción. La hipótesis de que $M'/xM'$ es íntegro (por tanto reducido)
implica entonces que $x$ es $M'$-regular en virtud de (3.4.6). Se concluye
que el núcleo de la homotecia $z\to xz$ en $M$ está contenido en
$\mathfrak p^jM\subset\mathfrak m^jM$ para todo entero $j$, y por tanto este
núcleo se reduce a $0$ $(\mathbf{0}_{\mathrm I},7.3.5)$, lo que demuestra que $x$ es
$M$-regular. Se puede entonces aplicar (3.4.5), lo que demuestra que $M$ es
íntegro.
\end{proof}

\begin{remark}[3.4.8]
\label{IV.2.3.4.8-fr}
La proposición análoga a (3.4.7), en la que se reemplaza «~íntegro~» por
«~reducido~», ya no es necesariamente cierta. Consideremos, por ejemplo, el
anillo de polinomios $C=K[X,Y,Z]$ sobre un cuerpo $K$, el anillo cociente
$B=C/\mathfrak p\mathfrak q$, donde $\mathfrak p=CZ$,
$\mathfrak q=CX^2+CY$; sea $A$ el anillo local de $B$ correspondiente al ideal
maximal imagen de $CX+CY+CZ$ en $B$. Si $x$, $z$ son las imágenes canónicas de
$X$, $Z$ en $A$, es claro que $xz\neq0$ pero $x^2z^2=0$; por otra parte, como
$A/xA$ es isomorfo a $K[Y,Z]/(YZ)$, se tiene $\dim(A/xA)=1$ mientras que
$\dim(A)=2$, por lo que $x$ pertenece a un sistema de parámetros de $A$
$(\mathbf{0},16.3.4)$, $A/xA$ es reducido, pero $A$ no lo es.
\end{remark}

\begin{proposition}[3.4.9]
\label{IV.2.3.4.9-fr}
Sean $A$ un anillo noetheriano, $M$ un $A$-módulo de tipo finito, $f$ un
elemento $M$-regular de $A$, tal que $M/fM$ no tiene ideales primos
asociados inmersos. Si $\mathfrak p_i$ $(1\leqslant\boldsymbol{i\leqslant m})$ son los
ideales primos asociados a $M/fM$, entonces, para todo entero $n>0$, $f^nM$
es la intersección de las imágenes inversas de $f^nM_{\mathfrak p_i}$ por
las aplicaciones canónicas $M\to M_{\mathfrak p_i}$ $(1\leqslant i\leqslant m)$.
\end{proposition}

\begin{proof}
Todo se reduce a demostrar que los $\mathfrak p_i$ son también los ideales
primos asociados a $M/f^nM$, pues entonces las saturaciones de los $f^nM$ respecto
a los $\mathfrak p_i$ son los submódulos de la descomposición primaria reducida
(necesariamente única) de $f^nM$ en $M$. Ahora bien, se tiene
\[
\operatorname{Ass}(f^{n-1}M/f^nM)
\subset\operatorname{Ass}(M/f^nM)
\subset\operatorname{Ass}(M/f^{n-1}M)
\cup\operatorname{Ass}(f^{n-1}M/f^nM)
\]
por (3.1.7), y como $f$ es $M$-regular,
$f^{n-1}M/f^nM$ es isomorfo a $M/fM$; basta entonces razonar por
inducción sobre $n$.
\end{proof}

\oldpage[IV-2]{52}
